\section{The clock carries the logits}
\label{sec:causal}

This is the section \citet{chen2026multiplication} name as missing. Their evidence for
stratified Fourier computation is, in their own words, \emph{``only correlational, meaning
causation cannot be directly inferred''}, and the future work they propose is exactly
\emph{``causal intervention techniques, such as ablations and activation patching''}. That is
why the ablation evidence here is a Results section rather than a line of support inside one.
\citet{nanda2023progress} and \citet{chughtai2023toy} fold ablations into their mechanism
sections, and we depart from that deliberately.

Sparsity says a representation is concentrated on a few characters. It does not say the
network computes with them. The protocol is Section~\ref{sec:methods-causal}'s: baseline, restricted
and excluded losses in exponent coordinates, against a null of $10{,}000$ random character
sets of the same cardinality, reported as a permutation $\hat{p}$.

\subsection{At prime powers}
\label{sec:causal-primepowers}

At five seeds across the six original moduli, not one of $10{,}000$ random equal-sized
character sets is as damaging as the key set in $26$ of the $31$ runs. The remaining five are
beaten by between $1$ and $2$ draws, so every run reads $\hat{p} \le 0.0003$:

\begin{center}
\begin{tabular}{llccc}
  \toprule
  $n$ & $(\mathbb{Z}/n\mathbb{Z})^{\times}$ & baseline/restricted, median & range over seeds & $p < 0.01$ in \\
  \midrule
  $113$ & $\mathbb{Z}_{112}$ & $2.34\times$ & $1.56$--$264$ & $5/5$ \\
  $121 = 11^{2}$ & $\mathbb{Z}_{110}$ & $1.85\times$ & $0.349$--$4.22$ & $5/5$ \\
  $125 = 5^{3}$ & $\mathbb{Z}_{100}$ & $1.98\times$ & $1.46$--$2.87$ & $5/5$ \\
  $119 = 7 \cdot 17$ & $\mathbb{Z}_{6} \times \mathbb{Z}_{16}$ & $0.67\times$ & $0.0146$--$1.96$ & $5/5$ \\
  $120$ & $\mathbb{Z}_{2}^{3} \times \mathbb{Z}_{4}$ & $1.25\times$ & $1.18$--$1.75$ & $5/5$ \\
  $165 = 3 \cdot 5 \cdot 11$ & $\mathbb{Z}_{2} \times \mathbb{Z}_{4} \times \mathbb{Z}_{10}$ & $1.42\times$ & $0.000709$--$1.63$ & $5/5$ \\
  \bottomrule
\end{tabular}
\end{center}

The improvement factor is reported as a median with its range across seeds, never as a mean.
Its distribution is heavy-tailed for the same reason the control distribution is. At $n = 113$
one seed of five reads $264\times$ where the other four read $1.56$ to $8.0$, and a mean over
those is $55.5 \pm 104.1$. That summary's standard deviation is twice its value, and one run
decides it. The claim rests on the permutation $p$, which is $5/5$
everywhere and is not a ratio at all.

Figure~\ref{fig:causal} gives the absolute losses and the null itself. At $n = 121$
seed $0$ the baseline loss is $4.96 \times 10^{-7}$, keeping only the key characters
\emph{improves} it to $2.10 \times 10^{-7}$, and removing exactly those characters raises it to
$1.29 \times 10^{1}$, seven orders of magnitude.
Keeping only the key characters improves the loss in $25$ of the $30$ five-seed runs, and in
$5/5$ seeds at $113$, $120$ and $125$.

The prime powers are the point. $121$ and $125$ are the local rings with nilpotents
and non-regular $\mathcal{J}$-classes, and their clock frequencies are load-bearing,
not merely a sparse description of the embedding.

\subsection{Where there is no discrete log}
\label{sec:causal-nodlog}

$119$, $120$ and $165$ have non-cyclic unit groups: no generator, no discrete log. They still
have an exponent coordinate, multiplication is still addition of exponent tuples, and a
character is a multi-index. Ablating there separates the key set from the median control by $10.0$ to
$4.6 \times 10^{7}\times$, so the
product-character circuit carries the logits too.

Across eleven further moduli at up to three seeds each, the permutation test passes at
$p < 0.01$ in $28$ of $29$ analysable runs, including $81 = 3^{4}$ at $3/3$. We name the
exception rather than the rate, and it is marginal rather than absent: $n = 98$ is $2/3$, with
one seed at $\hat{p} = 0.0151$ ($150$ of $10{,}000$ draws as damaging as the key set).

This count moved when the null was deepened, and the pre-registration said in advance that it
might. At $200$ draws it read $27$ of $29$. The second exception was $n = 49 = 7^{2}$ at
$p = 0.010$, two draws out of $200$, which landed exactly on the threshold.
\texttt{PREREGISTER\_r1\_permutation\_B.md} named both marginal tests before the re-run and
committed to reporting whatever they became, \emph{including} $26$ or $28$. At $10{,}000$
draws $n = 49$ reads $\hat{p} = 0.0050$ ($49$ draws) and passes. The other end of the range
did not move: $n = 98$'s seed read $0.015$ from $3$ of $200$ draws and reads $0.0151$ from
$150$ of $10{,}000$.

The pass at $n = 49$ is a caveat about the statistic, not a result about $n = 49$. That run
\emph{did not grok}: its held-out accuracy over the final ten samples is $0.891$, and its baseline loss, $1.3 \times 10^{-1}$, is four orders of magnitude
worse than any grokked run here. An earlier draft read its $p = 0.010$ as the test
``correctly'' reporting the absence of a circuit. It was not doing that, and the $50\times$
resolution is what exposed the reading. A permutation test on the logit spectrum asks whether
the key characters carry the logits; a model that has memorised much of its training grid
still has a logit tensor that is a function of $u \cdot v$ where it was trained. This is the
same limit Gate~2 measures directly. There $G1$, the same statistic and the named primary
criterion, passes on $20$ of $20$ \emph{failed} stratum-measurements
(Section~\ref{sec:strata-notclaimed}, Figure~\ref{fig:discrimination}). The permutation $p$
is evidence that a circuit is load-bearing \emph{where one exists}. It is not a detector of
whether the network generalised, and nothing in this paper asks it to be.

The analysis code was validated on a planted signal at $n = 35$ before it was trusted on
data.

\subsection{At $2^{7}$, the deepest nilpotency tested}
\label{sec:causal-128}

$n = 128$ is the sixth prime power and the hardest case in the set: $\omega = 1$, half its
residues zero divisors, unit group $\mathbb{Z}_{2} \times \mathbb{Z}_{32}$ with no discrete
log, and nilpotency depth $7$, the deepest tested. In product-character coordinates the key
characters carry the logits at the floor of the null, $\hat{p} = 1.0 \times 10^{-4}$, in $2/2$
analysable runs: $0$ of $10{,}000$ draws is as damaging as the key set in either.

The permutation $p$ is the claim here, and the improvement ratio is not. The ratio is given as
a median with its range. These samples are two or three runs drawn from a distribution
this paper has already shown to be heavy-tailed (Section~\ref{sec:methods-causal}), and a mean
over two or three such draws is exactly the summary we retired. Writing
baseline/restricted per run:

\begin{center}
\begin{tabular}{lcccc}
  \toprule
  $n$ & runs & median & range & $p < 0.01$ \\
  \midrule
  $128 = 2^{7}$        & $2$ & $239.52\times$ & $2.77$--$476.26$ & $2/2$ \\
  $131$                & $3$ & $2.88\times$   & $2.44$--$781.26$ & $3/3$ \\
  $127$                & $3$ & $5.04\times$   & $2.18$--$51.63$  & $3/3$ \\
  $91 = 7 \cdot 13$    & $2$ & $117.26\times$ & $1.45$--$233.07$ & $2/2$ \\
  $123 = 3 \cdot 41$   & $3$ & $1.76\times$   & $0.00$--$2.19$   & $3/3$ \\
  \bottomrule
\end{tabular}
\end{center}

\noindent
$13$ of $13$ runs across these five moduli reach the floor, $\hat{p} = 1/10{,}001 = 1.0 \times
10^{-4}$: no draw of $10{,}000$ was as damaging as the key set in any of them.
The ranges are the finding, and they are wide: at $n = 131$ one run of three improves by $781\times$ and the
other two by under $3\times$, so a mean of $262\times$ describes no run in the sample. An
earlier draft of this section reported exactly those means with standard deviations of the
same order as the values, which is the defect Section~\ref{sec:methods-causal} retires two
sections earlier. The ratios above supersede them. What survives unchanged is the
threshold-free statistic: every one of the $13$ runs is beyond every one of its $10{,}000$
draws.

Two qualifications. $n = 128$ grokked $2$ of $3$ seeds, so the $2/2$ is two analysable runs and
not two of three. And $n = 91$'s count is not counted as independent support. Our two analysis paths disagree
on its seed count, because one admitted a run at accuracy $0.976$. Under the three-state rule
that run is a near-grok and is neither data nor control.

The prime-power set therefore reaches $p = 2$ and depth $7$. The discrete-log
statistics of Section~\ref{sec:clock-primepowers} cover five cyclic prime powers; the ablation
covers all six.

\subsection{Independent of initialisation}
\label{sec:causal-init}

The dependence is not an artefact of one initialisation scale. Across $24$ runs at $n = 113$,
in three initialisation conditions (balanced fan-in, upstream-heavy and downstream-heavy,
following \citet{kunin2024getrich}) at eight seeds each, the permutation test passes at
$p < 0.01$ in $\mathbf{24/24}$. Every one of the $24$ is beyond \emph{every} one of its
$10{,}000$ draws.

\subsection{Inside the network, at a prime and at a prime square}
\label{sec:causal-internal}

Everything above acts on the logit tensor. It establishes what the output depends on, and
Section~\ref{sec:methods} states that it is not a statement about which internal
component computes it. At $n = 113$ and at $n = 121 = 11^{2}$ we can say the stronger thing.
The protocol, and the three criteria it inherits, are stated in
Section~\ref{sec:methods-causal}.

% STE descriptive
\begin{table}[t]
  \centering
  \small
  \begin{tabular}{cccccrl}
    \toprule
    $n$ & seed & baseline & restricted (acc.) & excluded (acc.) & excl./base. & $\hat{p}$ \\
    \midrule
    $113$ & $0$ & $1.2102\mathrm{e}{-05}$ & $2.0775\mathrm{e}{-06}$ ($1.0000$) & $4.7136\mathrm{e}{+01}$ ($0.0089$) & $3.895\mathrm{e}{+06}$ & $1/10001$ \\
    & $1$ & $1.2743\mathrm{e}{-07}$ & $3.2178\mathrm{e}{-07}$ ($1.0000$) & $3.5425\mathrm{e}{+01}$ ($0.0089$) & $2.780\mathrm{e}{+08}$ & $1/10001$ \\
    & $2$ & $1.5480\mathrm{e}{-07}$ & $2.2235\mathrm{e}{-07}$ ($1.0000$) & $2.3883\mathrm{e}{+01}$ ($0.0091$) & $1.543\mathrm{e}{+08}$ & $5/10001$ \\
    \midrule
    $121$ & $0$ & $1.7482\mathrm{e}{-06}$ & $3.4061\mathrm{e}{-07}$ ($1.0000$) & $5.0365\mathrm{e}{+01}$ ($0.0000$) & $2.881\mathrm{e}{+07}$ & $1/10001$ \\
    & $1$ & $7.7458\mathrm{e}{-03}$ & $1.0560\mathrm{e}{-01}$ ($0.9596$) & $5.7783\mathrm{e}{+01}$ ($0.0120$) & $7.460\mathrm{e}{+03}$ & $1/10001$ \\
    & $2$ & $1.6364\mathrm{e}{-03}$ & $2.2674\mathrm{e}{-03}$ ($1.0000$) & $1.8467\mathrm{e}{+01}$ ($0.0000$) & $1.128\mathrm{e}{+04}$ & $3/10001$ \\
    \bottomrule
  \end{tabular}
  \caption[Deletion of the key characters from the embedding returns the network to
    chance.]{\textbf{Deletion of the key characters from the embedding returns the network to
    chance.} Each row gives cross-entropy on the unit grid after an edit of $W_{E}$ and a new
    forward pass, three seeds per modulus. At $n = 113$ excluded accuracy $0.0089$--$0.0091$
    \emph{is} chance, $1/112 = 0.00893$. At $n = 121$ it is $0.0000$, $0.0120$ and $0.0000$
    against chance $1/110 = 0.00909$ on the $110 \times 110$ unit grid. Restricted accuracy is
    $1.0000$ on five of six runs and $0.9596$ on the sixth. The excluded/baseline median is
    $1.543 \times 10^{8}$ (range $3.895 \times 10^{6}$--$2.780 \times 10^{8}$) at
    $n = 113$.\par
    At $n = 121$ it is $1.128 \times 10^{4}$ (range $7.460 \times 10^{3}$--$2.881 \times
    10^{7}$). Both are medians and ranges over three seeds, because a standard deviation over
    three points is not a statistic. The two moduli are never pooled. $\hat{p}$ is
    Phipson--Smyth $(1+b)/(B+1)$ \citep{phipson2010permutation} against $B = 10{,}000$ random character sets of the same
    cardinality. Seeds $0$ and $1$ sit at the design floor and are printed as the fraction.
    Four decimals cannot distinguish ``at the floor'' from ``just above it''.}
  \label{tab:internal}
\end{table}
% STE end

At $n = 113$ all three criteria hold on all three seeds, against a pre-registered requirement
of two of three. The checkpoints are not new models. They reproduce the archived runs bit-identically,
$\max\lvert \Delta W_{E} \rvert = \max\lvert \Delta W_{U} \rvert
= \max\lvert \Delta \mathrm{logits} \rvert = 0$, with grokking steps $4{,}500$, $6{,}500$ and
$9{,}700$ recovered exactly. So Table~\ref{tab:internal} describes the same networks the rest
of this section measures. The projections were checked against the spectrum and not only
against the loss. Restricted preserves key-bin amplitude to the digit and zeroes the
complement to $10^{-13}$, and the key set holds $81.6\%$ of embedding amplitude in $8$ of $56$
bins at seed $0$.

\paragraph{The null's upper tail reaches the effect, and the permutation $p$ is what shows
it.} On seed $2$ the control \emph{maximum}, $2.4464 \times 10^{1}$, exceeds the excluded loss
$2.3883 \times 10^{1}$, and on seed $0$ it reaches $99.2\%$ of it. A ratio to the control
median would read $8.2 \times 10^{7}$ and hide this entirely. It is the same heavy-tail hazard
that retired the control-mean ratio, in a new arm. The cause is the design and it is
conservative: the draw pool includes the key characters, so a draw can partially reproduce
the intervention. Post hoc, and not part of the pre-registered analysis, the damage resolves into a
dose-response that is monotone \emph{in the mean} with how many key characters a draw removes.
Removals of $0$, $1$, $2$, $3$, $4$ and $5$ give $1.26 \times 10^{-5}$, $4.95$, $10.3$, $16.5$,
$23.9$ and $30.0$ at seed $0$. Figure~\ref{fig:internal} shows that the within-level spread is
wide, though, so the dose-response describes means and not individual draws. A draw containing \emph{no} key character does nothing at all, within $1.04$--$1.11
\times$ of baseline, and across the three seeds $12{,}468$ such draws never exceed
$1.05 \times 10^{-4}$ against an excluded loss of $24$--$47$. This is direct evidence for the
reading already given for the bimodality in Figure~\ref{fig:causal}: the harmless mode is the
draws the model does not use.

% STE descriptive
\begin{figure}[t]
  \centering
  \includegraphics[width=\textwidth]{../figures/F10_internal.pdf}
  \caption[The intervention acts on a weight, and the null's tail is draws that partially
    perform it.]{\textbf{The intervention acts on a weight, and the null's tail is draws that
    partially perform it.} \emph{Left:} baseline, restricted and excluded cross-entropy at
    $n = 113$ after an edit of $W_{E}$ and a new forward pass, three seeds. The panel uses the
    same three roles and hues as Figure~\ref{fig:causal}. Excluded is direct-labelled with its
    \emph{accuracy}, because that is the claim: $0.0089$--$0.0091$ is chance, $1/112$, not a
    degraded model. Restricted holds $1.0000$ on every seed.\par
    \emph{Right:} every one of the $30{,}000$ null draws, pooled over seeds and divided by its
    own run's baseline so the three share one axis. The draws are split by how many key
    characters each draw removed. Counts are above the axis, bars are medians, and the band is
    the real intervention, $3.895 \times 10^{6}$ to $2.780 \times 10^{8}$. The $12{,}468$ draws
    that remove none sit at the baseline line and never exceed $1.05 \times 10^{-4}$ in
    absolute terms. Damage rises monotonically in the mean with overlap, but the within-level
    spread is wide. A single key character can be worth four orders or almost nothing, so the
    dose-response is a statement about means, not about individual draws.\par
    This decomposition is post hoc and not part of the pre-registered analysis. The figure
    regenerates from \texttt{src/viz/plots.py::internal\_intervention}.}
  \label{fig:internal}
\end{figure}
% STE end

\paragraph{At a non-square-free modulus.} A prime is the weakest place to test this, because
every nonzero residue is a unit and the clause ``non-unit rows pass through untouched''
covers a single row. At $n = 121 = 11^{2}$ it covers eleven (zero and ten nilpotents),
and those rows form a non-regular $\mathcal{J}$-class at $d = 11$ (Section~\ref{sec:setup-strata}).
After the edit the network still holds eleven intact embedding rows and a whole non-regular
stratum. We ran the identical instrument, unmodified, under a separate pre-registration
committed before the retrain, with all three thresholds inherited from the $n = 113$ arm. The
pre-registration named $\text{excluded}/\text{baseline} < 100\times$ as the outcome that
would confine the internal claim to a prime. The worst of three seeds reads
$7.46 \times 10^{3}$, and all three criteria hold on all three seeds
(Table~\ref{tab:internal}). The checkpoints again reproduce the archived runs, with
$\max\lvert \Delta W_{E} \rvert = 0$ and grokking steps $10{,}700$, $10{,}300$ and $5{,}000$.

Three differences from $n = 113$ are reported, not interpreted. First, excluded accuracy falls
\emph{below} chance on two of three seeds ($0.0000$ against $1/110$) where at the prime it sat
on chance. No criterion reads accuracy against chance, so this is descriptive and would need
its own pre-registration to become a claim. Second, the null's tail again reaches the effect
on seed $2$. The control maximum $18.6887$ exceeds the excluded loss $18.4667$, and both
beating draws remove three of the four key characters. The $19{,}426$ draws that remove none
never exceed $1.71 \times 10^{-2}$ against an excluded loss of $18$--$58$. Third, the ratios
sit about two orders lower at the floor than at $n = 113$; the pre-registration declined to
predict that they would match, and we draw nothing from the comparison.

\paragraph{What this does and does not license.} At $n = 113$ and $n = 121$ the key
multiplicative characters are the representation the network \emph{computes with}, not merely
a basis its output is sparse in. That is the second of the two instruments
\citet{chen2026multiplication} name as missing, at two moduli and three seeds each, one of
them non-square-free with a non-regular stratum left intact. The logit-tensor result remains
the evidence at the other twenty-one moduli, and the scoping of Section~\ref{sec:methods}
stands there. Restricted is worse than baseline on two of three seeds at each modulus: by
$2.5\times$ and $1.4\times$ at $n = 113$, and by $13.6\times$ and $1.39\times$ at $n = 121$, at
accuracy $0.9596$ or above and at least two orders below excluded. This is the same
necessary-but-not-sufficient signature the non-cyclic case shows above, and a shortfall of the
key-set detector rather than of the circuit.

\subsection{Two caveats on scope}
\label{sec:causal-caveats}

Necessary is not the same as sufficient, and at two moduli we can only show the first. In the notation of \eqref{eq:ablation} the two are distinct claims:
\begin{equation}
  \label{eq:necsuf}
  \underbrace{\mathcal{L}(\Pi_{K^{c}} L) \;\gg\; \mathcal{L}(L)}_{K \text{ is \emph{necessary}}}
  \qquad\text{versus}\qquad
  \underbrace{\mathcal{L}(\Pi_{K} L) \;\lesssim\; \mathcal{L}(L)}_{K \text{ is \emph{sufficient}}} ,
\end{equation}
and only the first is what the permutation $p$ of \eqref{eq:permp} tests. At $n = 119$ the restricted loss is $2.0\times$ \emph{worse} than baseline: the six key
characters are necessary, since excluding them is catastrophic, but keeping only them does not
reconstruct the model. Both losses are around $10^{-5}$, six orders below the excluded loss,
so this is a shortfall in the frequency detector rather than a failed circuit, but it is not
the clean ``restricted improves'' signature the other runs show. $n = 123$ reproduces the same
pattern at a modulus that had never been seen ($p < 0.01$ in $3/3$, baseline/restricted median only
$1.76\times$, range $0.003$--$2.19\times$). The initialisation sweep shows it is condition-specific: $3/8$ under
upstream-heavy initialisation, $0/8$ under both of the others. It tracks the unbalanced
initialisation, not the modulus and not the group structure.

A frequency can sit in the embedding and do nothing. In our addition control at $n = 113$,
the frequency $k = 23$ has embedding norm $23.1$ against a noise floor near $1$, and a logit
ablation delta below $4 \times 10^{-10}$ (1 seed). It is present in
$W_{E}$ and absent from the circuit. \citet{nanda2023progress} report the same thing. The
consequence for method is that embedding-level frequencies are not circuit-level frequencies,
and the causal set is the one to report.

% STE descriptive
\begin{figure}[t]
  \centering
  \includegraphics[width=\textwidth]{../figures/F4_causal.pdf}
  \caption[The causal test and the null it rests on.]{\textbf{The causal test and the null
    it rests on.} \emph{Left:} absolute baseline, restricted and excluded cross-entropy per
    modulus, five seeds each. Bars are the median, and every seed is a dot. Removal of exactly
    the key characters costs three to seven orders of magnitude. A model with only them is at
    or below baseline. \emph{Right:} the permutation null for $n = 121$ seed $0$, $10{,}000$
    random character sets of equal cardinality, with the median control and the real key set
    marked.\par
    The null is bimodal. About three quarters of random removals leave the loss at baseline,
    because the model does not use those characters. The remainder cause damage but reach at most
    $10.9$. The key set is at $12.9$, beyond every one of the $10{,}000$ draws, so $\hat{p}$ is
    the floor, $1.0 \times 10^{-4}$. The bimodality is why a ratio to the control \emph{mean}
    is not a statistic here and was retired. The figure regenerates from
    \texttt{src/viz/plots.py::causal\_test}.}
  \label{fig:causal}
\end{figure}
% STE end
